% ============================================================================
% Working notes paper -- drafted from docs/working notes paper/backbone.txt
% and from the artifacts under final_system/ (SYSTEM.md, prf_rerank.py,
% final_system/system_submission/manifests/INDEX.md).
%
% PROVENANCE CONVENTION: every reported number carries an inline LaTeX comment
% stating exactly where it was read from, so it can be re-verified later
% without re-reading this whole file. "SYSTEM.md Sx" refers to the numbered
% section of final_system/SYSTEM.md (identical to the repository-root
% SYSTEM.md at the time of writing). "INDEX.md:<path>" refers to a row of
% final_system/system_submission/manifests/INDEX.md, keyed by the manifest
% path shown in that row. "prf_rerank.py:<def>" refers to a function in
% final_system/prf_rerank.py.
%
% PLACEHOLDER: Section 3.2 (web crawling) is intentionally left blank, as
% requested -- the crawler description will be supplied separately.
% ============================================================================
\documentclass[runningheads]{llncs}

% ================= PACKAGES & SETTINGS =================
\usepackage[utf8]{inputenc}
\usepackage[T1,T5]{fontenc}
\usepackage[vietnamese,english]{babel}
\usepackage{graphicx}
\usepackage{amsmath, amssymb}
\usepackage{booktabs}
\usepackage{hyperref}
\usepackage{float}
\usepackage{marvosym}

\begin{document}

% ================= TITLE & AUTHORS =================
% Working title; replace once the team confirms the final wording.
\title{A Multi-Run Hybrid Retrieval and Reranking System for Cross-Lingual Biomedical Information Retrieval}

\titlerunning{Multi-Run Hybrid Retrieval for Cross-Lingual Biomedical IR}

\author{
Author One\inst{1,2} \and
Author Two\inst{1,2} \and
Author Three\inst{1,2} \and
Corresponding Author\inst{1,2}(\Letter)
}
\authorrunning{Author One et al.}

\institute{
University of Information Technology, Ho Chi Minh City, Vietnam
\and
Vietnam National University, Ho Chi Minh City, Vietnam\\
\email{\{student1,student2,student3\}@gm.uit.edu.vn}\\
\email{\{advisor1,advisor2\}@uit.edu.vn}
}

\maketitle

% ================= ABSTRACT =================
\begin{abstract}
Retrieving biomedical evidence for a query written in a low-resource language,
from a document collection that spans several languages, poses two coupled
difficulties: bridging the language gap between the query and a large share of
the collection, and deciding, for each query independently, how many documents
and passages are worth returning. We describe a system built for this setting
that is assembled entirely from openly released components, with no
task-specific fine-tuning. First, a multi-run hybrid retrieval stage combines
sparse (BM25) and dense (bi-encoder) retrieval through reciprocal rank fusion
across several independently prepared document partitions, one of which
retrieves against a machine-translated query to reach documents in the
collection's dominant language. Second, a higher-capacity cross-encoder
reranker rescores the union of candidates surfaced by every partition. Third,
a pseudo-relevance feedback step re-ranks documents using a per-query topic
vector estimated from the current top results, and the number of documents
returned per query is set by directly maximizing an expected-$F_2$ objective
under a probabilistic reading of the reranker scores, rather than by a fixed
cutoff. Passages are then cut as token-bounded windows around the
highest-scoring span of each selected document, so that every submitted
passage remains a verbatim, contiguous excerpt of the source text. On the
public leaderboard of the task this pipeline, together with a coverage-recovery
step that retrieves archived snapshots of pages the crawler could not fetch
directly, reaches a final score of 0.3060 (document-level $F_2$ 0.3381,
passage-level $F_2$ 0.2739), with each refinement in the pipeline contributing
a measurable, monotonic gain over the preceding configuration. We report this
progression together with several design alternatives that were tried and
found not to improve the score, and we discuss the main limitations of
evaluating exclusively against a rate-limited public leaderboard.

\keywords{Cross-lingual information retrieval \and Biomedical text retrieval \and Hybrid dense-sparse retrieval \and Reranking \and Pseudo-relevance feedback \and Adaptive result-set sizing}
\end{abstract}

% ================= 1. INTRODUCTION =================
\section{Introduction}

Access to biomedical evidence is unevenly distributed across languages: a
substantial share of consumer-facing medical content for languages such as
Vietnamese exists only as translations or summaries of material originally
published in English or Chinese, while the underlying primary sources remain
in their original language. A retrieval system that is asked to answer a
query written in one language from a collection that mixes several languages
must therefore locate relevant evidence regardless of which language a
particular document happens to be written in, rather than assuming the query
and the answer share a language. This setting, together with the practical
constraint that only openly released models may be used and that no
model-specific labelled training data is provided, motivates building a
retrieval pipeline out of general-purpose, pretrained components rather than
a system trained end-to-end for this task.

A first difficulty is that lexical (sparse) retrieval cannot bridge a
vocabulary mismatch across scripts: a Vietnamese query and a Chinese document
share almost no surface tokens, so a system that relies on term overlap alone
systematically under-retrieves the language in which most of the collection is
written~\cite{lipton2014thresholding}. Dense bi-encoder retrieval partially
closes this gap when the encoder was trained on multilingual data, but recent
reports on multilingual embedding and reranking models still show a
non-trivial spread in cross-lingual effectiveness depending on the script
pair and the domain~\cite{qwen3embedding2025}, which suggests that no single
encoder or reranker should be trusted as sufficient on its own for a
collection this heterogeneous.

A second difficulty is specific to the evaluation protocol used here: the
number of relevant documents and passages is not the same for every query, yet
the official metric rewards returning close to the right number of items
rather than either a fixed short list or an exhaustive long list. A system
that always returns, say, the top 50 documents for every query will be
unnecessarily imprecise on queries with few relevant answers and will miss
recall on queries with many. Treating the per-query result-set size as a
decision variable, rather than a global constant, requires a principled way
to trade expected precision against expected recall using only the scores the
system itself produces, since no labelled development set matching the
official judgments is available.

We address these two problems through three component-level contributions
that compose into a single pipeline. First, a \emph{multi-run hybrid retrieval}
stage (Sect.~\ref{sec:module1}) runs sparse-dense fusion independently over
several partitions of the crawled corpus -- including one partition retrieved
against a machine-translated form of the query -- so that recall is not
limited by any single encoder, tokenizer, or query language. Second, a
\emph{higher-capacity reranker rescoring} stage (Sect.~\ref{sec:module2})
consolidates the candidates surfaced by every partition and rescores them with
a larger cross-encoder than the one used during first-stage retrieval,
spending the more expensive model only on the comparatively small candidate
pool rather than on the full collection. Third, a \emph{pseudo-relevance
feedback and adaptive selection} stage (Sect.~\ref{sec:fusion}) re-ranks the
rescored candidates using a per-query topic estimate, and chooses how many
documents and passages to submit for each query by directly optimizing an
expected-$F_2$ criterion derived from the reranker scores, subject to
practical bounds and to an overall submission size limit.

We evaluate the complete pipeline, together with each of the above
refinements added one at a time, on the public leaderboard of the shared task
that supplied the queries, the document identifiers, and the official metric
(Sect.~\ref{sec:dataset}). Starting from a pseudo-relevance-feedback baseline
that reorders a fixed-size candidate pool, each subsequent refinement that we
report keeping -- per-query adaptive sizing, per-query passage windowing, a
second passage window for highly-ranked documents, and a supplementary round
of reranker scoring near the per-query cutoff -- is associated with a
further increase in the leaderboard score, reaching 0.3060 overall
(Sect.~\ref{sec:results}). We also report the alternatives we tried and did
not keep, because they left the score unchanged or lowered it, since we
believe that a record of what did not work is as useful to future
replications as a record of what did.

% ================= 2. RELATED WORK =================
\section{Related Work}

\subsection{Cross-Lingual and Multilingual Biomedical Retrieval}
Cross-lingual information retrieval has traditionally relied on either
translating the query into the document language, translating documents into
the query language, or mapping both into a shared multilingual representation
space~\cite{clirudit2025}. Query translation avoids the cost of translating an
entire collection but is reported to perform less reliably than document-side
translation or than translation-free dense retrieval, in part because short
queries carry less context for a translation model to resolve
ambiguity~\cite{clirudit2025}. Dedicated multilingual machine translation
models aimed at this use case, such as those built specifically to cover a
Vietnamese-Chinese-English language mix, have recently been
released~\cite{hunyuanmt2025}; our pipeline uses general-purpose query
translation rather than such a dedicated system, which we view as an
available avenue for improvement rather than as a closed question. Benchmark
efforts aimed specifically at Vietnamese text embeddings, built by machine
translating existing multilingual evaluation suites, have also started to
appear~\cite{vnmteb2025}, reflecting that cross-lingual retrieval quality for
Vietnamese specifically is still an actively measured, rather than solved,
problem.

\subsection{Dense and Hybrid Retrieval with Learned Rerankers}
Reciprocal rank fusion combines several ranked lists by the reciprocal of
each item's rank rather than by its raw score, which avoids having to
calibrate heterogeneous score scales -- a lexical BM25 score and a dense
cosine similarity, in particular -- onto one another before they can be
combined~\cite{cormack2009rrf}. A two-stage retrieve-then-rerank design, in
which an inexpensive first-stage retriever proposes a candidate pool that a
more expensive cross-encoder then rescores, is standard practice when the
reranker's per-pair cost makes scoring the full collection impractical.
General-purpose multilingual embedding and reranking model families released
in various parameter sizes, evaluated on standard retrieval and reranking
benchmarks, are reported to perform competitively across several languages
and model sizes~\cite{qwen3embedding2025}; our pipeline draws its dense
encoders and cross-encoder rerankers from this kind of openly released,
general-purpose family, used here without any task-specific fine-tuning.

\subsection{Pseudo-Relevance Feedback and Adaptive Result-Set Sizing}
Pseudo-relevance feedback re-ranks or re-weights retrieval results using
information drawn from the system's own top-ranked results, under the
assumption that most of them are on-topic even without external relevance
judgments; we use this idea in its simplest dense form, representing the
``topic'' of a query as the centroid of the embeddings of its currently
top-ranked passages. Separately from how documents are ranked, deciding how
many of them to return is a thresholding problem: for a classifier whose
output can be read as a calibrated probability of relevance, the decision
threshold that maximizes the expected $F_1$ score is known to equal half of
the $F_1$ value attained at that threshold, and the corresponding result for
$F_\beta$ follows the same derivation~\cite{lipton2014thresholding}. Our
adaptive document-count rule (Sect.~\ref{sec:fusion}) is a direct,
per-query application of that $F_\beta$-optimal thresholding argument to the
reranker's output scores, rather than a fixed cutoff tuned once on a held-out
set, which this task does not provide.

% ================= 3. PROPOSED SYSTEM =================
\section{Proposed System}

\subsection{System Overview}
\label{sec:overview}

% PLACEHOLDER (figure to be supplied together with the crawler section).
\begin{figure}[!htbp]
    \centering
    \includegraphics[width=\textwidth]{figs/fig-system-overview.pdf}
    \caption{Overview of the proposed system: an offline corpus-construction
    phase (web crawling and archived-snapshot recovery) feeds a query-time
    ranking phase (multi-run hybrid retrieval, reranker rescoring, and
    pseudo-relevance-feedback--based adaptive selection).}
    \label{fig:system}
\end{figure}

Our system consists of a corpus-construction phase and a query-time ranking
phase (Fig.~\ref{fig:system}). The corpus-construction phase
(Sect.~\ref{sec:crawl}) turns the organizers' list of source URLs into
locally held, quality-checked document text; because not every listed URL
could be fetched directly, a coverage-recovery step
(Sect.~\ref{sec:coverage}) separately resolves a subset of the remaining URLs
through their archived snapshots. The query-time ranking phase processes each
of the 1{,}200 queries through three stages: multi-run hybrid retrieval
(Sect.~\ref{sec:module1}), higher-capacity reranker rescoring
(Sect.~\ref{sec:module2}), and pseudo-relevance feedback with adaptive
document and passage selection (Sect.~\ref{sec:fusion}). The output of the
last stage is serialized directly into the submission format required by the
task, with every passage traced back, by character offset, to an exact,
contiguous span of the corresponding document's crawled text, so that no
submitted passage is paraphrased or regenerated.
% PROVENANCE: the exact-span / non-rewritten guarantee is implemented in
% cut_chunks() in final_system/prf_rerank.py (slices text[start:end] from the
% run's corpus) and is stated as a requirement in SYSTEM.md S3 ("chunk_text
% luôn là một lát cắt liên tục, nguyên văn của văn bản tài liệu").

\subsection{Corpus Construction via Web Crawling}
\label{sec:crawl}

% PLACEHOLDER: this subsection is intentionally left blank. The description
% of the crawling pipeline (source list, crawler design, per-run partitioning,
% quality control, and the resulting per-run document counts) will be
% supplied separately and inserted here.
[To be completed.]

\subsection{Coverage Recovery via Archived Snapshots}
\label{sec:coverage}

Not every URL in the organizers' source list could be retrieved through a
direct request at crawl time. For the documents that remained missing after
the main crawl, we recover a subset of them from publicly archived snapshots
of the same URL, under the same identifier, rather than discarding the query
coverage those URLs would otherwise contribute.
% PROVENANCE: SYSTEM.md S4 "B2. Tìm trang chưa crawl được theo URL" and the
% tool list (slug_retrieval.py, slug_rerank.py, slug_targets.py,
% final_system/notebooks/archive_fetch_kaggle.ipynb).

This recovery proceeds in four steps. First, because the URL itself is the
only signal available for a page that was never fetched, we build a sparse
(BM25) index over the tokens of each unfetched URL's slug and retrieve
candidate URLs for every query against that index. Second, a general-purpose
language model scores each (query, candidate) pair using the slug-derived
title as a proxy for the page's topic, which lets us rank candidate URLs by
plausible relevance before spending any crawling effort on them. Third, the
highest-scoring candidate URLs are compiled into a target list. Fourth, for
each targeted URL we query the Wayback Machine for an archived snapshot and,
when one is available, treat its content as the document text for that
identifier, subject to the same quality control as the main crawl.
Provenance for every archived page -- the archived timestamp and the
originating URL -- is retained alongside the recovered text.
% PROVENANCE: SYSTEM.md S3 "Dữ liệu ngoài (web.archive.org)"; provenance
% files under work/archive*/ (targets.parquet + manifest), listed in
% final_system/system_submission/manifests/INDEX.md under archive/,
% archive2/, archive3/.

Documents recovered this way are processed by the same per-run pipeline
described in Sect.~\ref{sec:module1} as three additional partitions
(denoted \texttt{run\_archive}, \texttt{run\_archive2}, and
\texttt{run\_archive3} below), rather than being merged into the partitions
built from directly crawled text, so that their provenance remains
separately auditable.

\subsection{Multi-Run Hybrid Retrieval}
\label{sec:module1}

The crawled and recovered documents are not processed as one monolithic
corpus. They are instead split into several independently prepared
\emph{runs} -- contiguous batches of source identifiers, crawled and indexed
separately -- and queried independently; results are combined only after
retrieval. This design keeps indexing incremental (a newly crawled batch
becomes its own run rather than requiring the full index to be rebuilt) and
lets us vary, per run, which query form is used to search it.
% PROVENANCE: SYSTEM.md S4 ("B1. Pipeline từng run") lists the 14 runs kept
% in the final system: run_vi, run_b..run_g, run_zh_a/b/c, run_h, run_archive,
% run_archive2, run_archive3.

Within each run, a document is split into overlapping passages of up to
$320$ tokens, with an overlap of $80$ tokens and a hard ceiling of $768$
tokens per passage, using sentence and paragraph boundaries where available;
each passage retains the character offsets of its source span. Each run is
indexed twice: a sparse BM25 index over the passage text, and a dense index
built with a general-purpose multilingual bi-encoder. Retrieval combines a
sparse channel, a dense channel using \texttt{BAAI/bge-m3}, and a second
dense channel using \texttt{Qwen/Qwen3-Embedding-0.6B}, fused by reciprocal
rank fusion (RRF)~\cite{cormack2009rrf} rather than by combining raw scores,
precisely because BM25 scores and cosine similarities are not on a comparable
scale. The fused candidates for each query are then rescored by a first-stage
cross-encoder reranker -- \texttt{BAAI/bge-reranker-v2-m3} and
\texttt{Qwen/Qwen3-Reranker-0.6B}, run in sequence within the same run -- to
produce the per-run candidate score that feeds Sect.~\ref{sec:module2}.
% PROVENANCE: SYSTEM.md S4 "B1" bullet list and S2 model table; chunking
% parameters (320/80/768) match src/vibiomir/pipeline.py defaults referenced
% throughout docs/solution/run_guide.md and training_design.md.

One run, denoted \texttt{run\_zh\_a}, departs from this template in its query
side only: because a majority of the crawled collection is written in
Chinese (see Sect.~\ref{sec:dataset}), the original Vietnamese query is first
translated into Chinese by a general-purpose instruction-tuned language model,
\texttt{Qwen/Qwen3-4B} in 4-bit quantization, and the translated form is used
to search that run's sparse and dense indexes. This is the only point in the
pipeline at which a generative language model produces text that influences
retrieval; it never produces text that is itself submitted as evidence.
% PROVENANCE: SYSTEM.md S2 model table (Qwen/Qwen3-4B row, "Dịch câu hỏi
% sang tiếng Trung cho run_zh_a") and S4 "run_zh_a dịch câu hỏi sang tiếng
% Trung bằng Qwen3-4B."

\subsection{Higher-Capacity Reranker Rescoring}
\label{sec:module2}

The first-stage rerankers in Sect.~\ref{sec:module1} are chosen to be cheap
enough to run within each of the fourteen partitions separately. We treat
their output only as a way to assemble a manageable, union candidate pool
across runs, and reserve a substantially larger cross-encoder,
\texttt{Qwen/Qwen3-Reranker-4B}, to rescore that pool once, after the
per-run candidates have been merged by query. For every (query, document)
pair in the pool, we score up to three of the document's highest-ranked
passages and keep, for that pair, the highest of those three scores together
with the identity of the scored passage, since the most informative passage
of a document is not necessarily its first-stage-highest-ranked one.
% PROVENANCE: SYSTEM.md S4 "B3" ("Mỗi tài liệu ứng viên được chấm 3 đoạn tốt
% nhất") and S6 pairs_manifest.json rows listing PASSAGE_RANKS=3 behaviour
% (pool_passages() docstring in final_system/prf_rerank.py).

Because this stage uses a markedly larger model than the rest of the
pipeline, it is run separately, on paired GPU accelerators, rather than as
part of the per-run pipeline; the candidate pool is exported once per pooling
round and scored in that separate environment. We used five pooling rounds
over the course of development, each covering a different stage of the
candidate pool as additional runs (including the archive-recovered runs of
Sect.~\ref{sec:coverage}) became available, plus one supplementary round that
rescores only the passages of documents sitting near each query's eventual
document-count cutoff (Sect.~\ref{sec:fusion}) that had not yet been scored
by any earlier round, so that the adaptive cutoff decision is not made on an
incompletely scored neighbourhood of candidates.
% PROVENANCE: SYSTEM.md S4 "B3" round list (rerank_pool4, rerank_pool5,
% rerank_h_arch, rerank_arch2, rerank_arch3, rerank_extra) and S5 ("Lệnh tái
% lập") giving pairs_total=2{,}067{,}495 after all rounds and 165{,}331 pairs
% for the supplementary round; export_extra() / f2_counts() in
% final_system/prf_rerank.py implement the near-cutoff selection.

\subsection{Pseudo-Relevance Feedback, Fusion, and Adaptive Selection}
\label{sec:fusion}

The last stage of the pipeline turns the rescored candidate pool into the
submitted ranking and result-set size for each query, and is implemented
without any further GPU inference. It combines four decisions: how to fuse
the two reranker scores available for a document, how to re-rank documents
using a query-level topic estimate, how many documents and passages to
return for a given query, and how to cut a submitted passage.

\paragraph{Document scoring.}
Each candidate document already carries two cross-encoder scores: the
first-stage score $s_{0.6B}$ from Sect.~\ref{sec:module1} and the
higher-capacity score $s_{4B}$ from Sect.~\ref{sec:module2}, when the latter
is available. We combine them linearly,
\begin{equation}
s_{\mathrm{doc}} = \mu \cdot s_{4B} + (1-\mu) \cdot s_{0.6B}, \qquad \mu = 0.5,
\label{eq:mix}
\end{equation}
so that a document that was pooled but has not yet received a $4$B score
degrades gracefully to its first-stage score rather than being dropped.
% PROVENANCE: pool_passages() in final_system/prf_rerank.py, default
% --mix 0.5 (argparse default in main()); SYSTEM.md S4 step B4.1.

\paragraph{Pseudo-relevance feedback.}
For each query, we form a topic vector as the mean \texttt{bge-m3} embedding
of the best-scoring passage of the top $10$ documents under
Eq.~\eqref{eq:mix}, then re-rank the candidate pool by
\begin{equation}
s_{\mathrm{final}} = z(s_{\mathrm{doc}}) + w \cdot z\!\left(\max_{c \in \mathrm{chunks}(d)} \cos(c, \mathbf{t})\right), \qquad w = 1,
\label{eq:prf}
\end{equation}
where $z(\cdot)$ is the sample $z$-score over the candidate pool of a given
query, $\mathbf{t}$ is the query's topic vector, and the maximum is taken
over the passage embeddings of document $d$. Standardizing both terms before
summing them avoids letting whichever term happens to have the larger raw
scale dominate the ranking.
% PROVENANCE: zscore(), and the "base + args.weight * topic" combination in
% the main orchestration of final_system/prf_rerank.py; SYSTEM.md S4
% step B4.2 (weight $1 \times$ cosine of best chunk against the topic
% vector of the top 10 documents).

\paragraph{Adaptive document count.}
Rather than returning a fixed number of documents for every query, we choose
a per-query count $N_q$ that maximizes an expected-$F_2$ objective. Reading
the (calibrated, rescaled) reranker score as an estimate of the probability
that a candidate is relevant, let $S_N = \sum_{i=1}^{N} v_i$ be the
cumulative sum of the top-$N$ scores under Eq.~\eqref{eq:prf}'s ranking, and
let $G_q = c \cdot T_q / \rho$ be the estimated number of relevant documents
for query $q$, where $T_q$ is the score total over the whole candidate pool,
$\rho$ is the fraction of official relevant documents estimated to fall
inside the candidate pool (measured empirically at $\rho \approx 0.73$), and
$c$ is a single global calibration constant. The per-query count is then
\begin{equation}
N_q = \operatorname*{arg\,max}_{\,\mathrm{low} \,\leq\, N \,\leq\, \mathrm{high}} \ \frac{c \cdot S_N}{4\, G_q + N},
\label{eq:f2}
\end{equation}
which is, up to the constant factor $5$, the closed-form optimum of
$F_2 = 5PR/(4P+R)$ under $P \approx c S_N / N$ and $R \approx c S_N / G_q$ --
the per-query analogue, under this probabilistic reading of the scores, of
the result that the $F_\beta$-optimal decision threshold on a calibrated
score is a fixed function of the achieved $F_\beta$
value~\cite{lipton2014thresholding}. The single constant $c$ is calibrated,
by bisection, so that the mean of $N_q$ over all queries equals a target value
(190 in the configuration reported in Sect.~\ref{sec:results}), and $N_q$ is
clipped to a plausible range (60--400 documents) to avoid degenerate counts
on queries whose score distribution is unusually flat or peaked.
% PROVENANCE: f2_counts() docstring and implementation in
% final_system/prf_rerank.py ("F2(N) = 5·c·S_N/(4·G+N)", pool_recall default
% and bisection over log(c)); SYSTEM.md S4 step B4.3 (bounds 60--400, mean
% target 190).

\paragraph{Passage windowing and submission size.}
For every selected document, we submit one window of up to $1{,}000$ tokens
built from whole sentence/paragraph segments around the document's
highest-scoring passage under Sect.~\ref{sec:module2}; the window is still an
exact, contiguous slice of the source text, never a concatenation of
non-adjacent spans. For the $50$ highest-ranked documents of each query, a
second window is added around the second-highest-scoring passage of the same
document whenever that passage falls outside the first window, so that a
long document with two separately relevant sections is not forced to
compete for a single window. The resulting submission JSON is compressed
with LZMA at a $512$~MB dictionary size, which we verified byte-for-byte
against the uncompressed JSON before using it, to keep the submitted archive
under the task's size limit while introducing no further rewriting of the
submitted text.
% PROVENANCE: SYSTEM.md S4 step B4.4-B4.5 (window $\le 1{,}000$ tokens,
% second window for top 50 documents, JSON capped near 1{,}050 MB) and S5
% (rezip_lzma.py, --dict-mb 512, byte-level equality check); cut_chunks() in
% final_system/prf_rerank.py for the verbatim-slice guarantee.

% ================= 4. EXPERIMENTS =================
\section{Experiments}

\subsection{Dataset and Evaluation Setup}
\label{sec:dataset}

The task supplies $1{,}200$ Vietnamese-language biomedical queries and a
fixed list of source identifiers and URLs; participants are responsible for
crawling the corresponding documents themselves, and only documents on this
official identifier list may be submitted as evidence. Summed over the
fourteen retained runs of Sect.~\ref{sec:module1}--\ref{sec:coverage}, the
per-run document counts reported by the corpus-preparation step total
$796{,}579$ kept documents and $2{,}763{,}929$ passages before any
cross-run deduplication; because the same underlying document can in
principle be crawled into more than one run, this sum should be read as an
upper bound on the number of distinct documents available to the system,
not as a deduplicated corpus size.
% PROVENANCE: sum of documents_kept / chunks over the 14 runs listed in
% SYSTEM.md S4 (run_vi, run_b..run_g, run_zh_a/b/c, run_h, run_archive,
% run_archive2, run_archive3), each read from the corresponding
% */corpus/prepare_manifest.json and */chunks/chunk_manifest.json row in
% final_system/system_submission/manifests/INDEX.md. Three additional runs
% present in that index (run_en_a, run_en2, run_en3, together 239{,}441
% documents) correspond to a PubMed-based English expansion that SYSTEM.md
% S3 records as tried and not kept ("Không dùng: PubMed (đã thử, không có
% lợi)"); they are excluded from this total.

Performance is measured by the organizers at two levels of granularity,
document and passage, using macro-averaged Precision, Recall, and $F_2$ over
all $1{,}200$ queries; a submitted passage counts as correct for a given
reference passage when the longest common token subsequence between the two,
after a fixed normalization, covers at least $40\%$ of the reference
passage's length. We report the aggregate FINAL score published by the
leaderboard together with the document-level and passage-level $F_2$ scores
that compose it. All results below are measured on the public leaderboard of
the shared task; the task additionally defines a separate private evaluation
phase over the same submission format, which this paper does not report on
separately. No labelled development set matching the official judgments was
available to us during development; the ablations in
Sect.~\ref{sec:ablation} were therefore evaluated directly against the public
leaderboard, a methodological limitation we return to in
Sect.~\ref{sec:limitations}.
% PROVENANCE: query count, document-identifier policy and the 40% LCS
% passage-matching rule are stated in docs/note/RULE.txt (sections 3.1-3.2
% of the task rules); FINAL/DOCS_F2/CHUNKS_F2 column names and values are
% read from SYSTEM.md S1 and S7.

\subsection{System-Level Results}
\label{sec:results}

Table~\ref{tab:main} reports the two configurations we submitted as final
candidates. Both apply every refinement described in
Sect.~\ref{sec:fusion}; the only difference between them is the passage
compression budget used by the fitting step that trades the number of
second-window passages against the JSON size limit (Sect.~\ref{sec:fusion},
``Passage windowing and submission size''). We designated the slightly
higher-scoring configuration as our primary submission.

\begin{table}[H]
\centering
\caption{Public-leaderboard scores of the two final submission candidates,
over $1{,}200$ queries.}
\label{tab:main}
\footnotesize
\setlength{\tabcolsep}{4pt}
\begin{tabular}{lcccc}
\hline
Configuration & FINAL & DOCS\_F2 & CHUNKS\_F2 & Role \\
\hline
\texttt{f2\_x50\_rx\_small} & \textbf{0.3060} & \textbf{0.3381} & \textbf{0.2739} & primary \\
\texttt{f2\_x50\_1050\_small} & 0.3055 & 0.3375 & 0.2734 & backup \\
\hline
\end{tabular}
\end{table}
% PROVENANCE: SYSTEM.md S1 table, verbatim.

\subsection{Incremental Contribution of Each Refinement}
\label{sec:ablation}

We tracked the public-leaderboard score as each refinement described in
Sect.~\ref{sec:fusion} was added on top of the previous configuration, keeping
everything else fixed. Table~\ref{tab:progression} reports this progression;
each row's FINAL score already includes every refinement in the rows above
it. The step contributed by the adaptive, per-query document count is
reported for the document-level score only, because that intermediate
configuration was submitted before passage windowing had been added to the
pipeline, and we did not re-run a since-superseded configuration purely to
backfill a passage-level number.

\begin{table}[H]
\centering
\caption{Public-leaderboard FINAL score after each refinement is added, in
the order it was introduced. Each row includes all refinements listed above
it.}
\label{tab:progression}
\small
\setlength{\tabcolsep}{4pt}
\begin{tabular}{lccl}
\hline
Refinement & FINAL & DOCS\_F2 & Note \\
\hline
PRF re-ranking + chunk cutting by reranker order & 0.2851 & 0.3158 & fixed-size candidate pool \\
+ Per-query adaptive document count & n/a$^{*}$ & 0.3336 & document-level only \\
+ Per-query passage windowing & 0.3003 & --\, & CHUNKS\_F2 0.2669 \\
+ Expected-$F_2$ count rule, 1{,}000~MB budget & 0.3033 & 0.3375 & \\
+ Second window, 20 $\rightarrow$ 50 documents & 0.3050 $\rightarrow$ 0.3055 & --\, & CHUNKS\_F2 0.2725 $\rightarrow$ 0.2734 \\
+ Supplementary near-cutoff rescoring & \textbf{0.3060} & 0.3381 & CHUNKS\_F2 0.2739 \\
\hline
\end{tabular}

\smallskip
\footnotesize $^{*}$Only the document-level submission of this intermediate
configuration was scored; a FINAL score for a document-only submission is
not comparable to the full-pipeline rows above and is omitted here.
\end{table}
% PROVENANCE: SYSTEM.md S7, "Các bước đã thử và đóng góp" table, verbatim
% (including the footnote content, paraphrased from "ước tính, chỉ nộp
% phần docs").

Every refinement in Table~\ref{tab:progression} is associated with a FINAL
score at least as high as the configuration before it, and the overall gain
from the first row to the last is $0.0209$ FINAL. We read this as evidence
that the three mechanisms of Sect.~\ref{sec:fusion} -- topic-based
re-ranking, per-query adaptive sizing, and passage windowing -- address
largely separate sources of error rather than one mechanism substituting for
another, since each continued to help after the others were already in
place.

\subsection{Alternatives That Were Not Kept}
\label{sec:rejected}

We also evaluated several modifications that either left the score
essentially unchanged or reduced it; Table~\ref{tab:rejected} lists them
together with the document-level $F_2$ obtained, measured against the
0.3336 or 0.3375 document-level baselines of the corresponding row in
Table~\ref{tab:progression} (the two reference points we had available at
the time each alternative was tried).

\begin{table}[H]
\centering
\caption{Design alternatives evaluated and not retained, with the resulting
document-level $F_2$ on the public leaderboard.}
\label{tab:rejected}
\footnotesize
\setlength{\tabcolsep}{4pt}
\begin{tabular}{lc}
\hline
Alternative & DOCS\_F2 \\
\hline
Expanding candidates beyond the pool by topic similarity & 0.3121 \\
Qwen-based encoder for the PRF topic vector & 0.3327 \\
Mean target of 230 documents per query & 0.3324 \\
Maximum of 600 documents per query & 0.3306 \\
Minimum of 30 documents per query & 0.3324 \\
Down-weighting Chinese-language documents by $\times 0.6$ & 0.3294 \\
Folding the second passage's score into the document score & 0.3360 \\
Mixing coefficient $\mu = 0.6$ in Eq.~\eqref{eq:mix} & lower \\
Minimum of 90 documents per query & lower \\
\hline
\end{tabular}
\end{table}
% PROVENANCE: SYSTEM.md S7, "Các thử nghiệm bị loại" table, verbatim (two
% rows reported only as "thấp hơn" in the source, reproduced as "lower"
% since no exact value was recorded there).

We note two alternatives that a reader might expect to help but did not:
expanding the candidate pool with topically similar but previously unranked
documents reduced the score rather than increasing it, and down-weighting
Chinese-language documents -- motivated by the Vietnamese language of the
queries -- also reduced the score. We read the latter result as consistent
with the corpus composition described in Sect.~\ref{sec:dataset}: because a
large share of the collection is Chinese-language, suppressing it is likely
to remove genuinely relevant documents more often than it removes
false positives, even though the query language itself is Vietnamese.

\subsection{Limitations}
\label{sec:limitations}

We highlight four limitations of the present system and of how it was
evaluated. First, the ablations reported in Sect.~\ref{sec:ablation} and
Sect.~\ref{sec:rejected} were compared directly on the public leaderboard
rather than on an independent, labelled development set, because no such set
matching the official judgments was available to us; the task additionally
limits the number of submissions allowed per day, which constrains how many
configurations can be compared this way and raises the ordinary risk of
indirectly fitting the public split through repeated, score-guided
decisions. Second, the reported corpus size of $796{,}579$ documents is a
sum of per-run counts, not a deduplicated figure, since no explicit
cross-run deduplication step is applied before the per-run statistics are
computed; the number of distinct documents actually available to the system
is therefore somewhat smaller than this figure. Third, the coverage-recovery
step of Sect.~\ref{sec:coverage} depends on a URL's slug carrying enough
signal for the first retrieval sub-step to rank it, and on that URL having
been archived publicly, so it is expected to recover only a fraction of the
documents the main crawl could not fetch rather than closing the coverage
gap entirely. Fourth, every learned component in the pipeline is used in a
zero-shot, pretrained configuration; we did not fine-tune any encoder or
reranker on this task's queries or documents, so headroom that a
task-specific fine-tuning step might provide is not reflected in the results
above.
% PROVENANCE: 796,579 figure is this paper's own sum, reproduced from the
% per-run manifest table under Sect.~\ref{sec:dataset}'s provenance note,
% not an independently stated total in SYSTEM.md.

% ================= 5. CONCLUSION =================
\section{Conclusion}

We presented a retrieval pipeline for cross-lingual biomedical information
retrieval that is assembled from openly released, pretrained components --
sparse and dense first-stage retrieval fused by reciprocal rank fusion across
several independently maintained document partitions, a higher-capacity
cross-encoder reranker applied only to the resulting candidate pool, and a
pseudo-relevance-feedback stage that both re-ranks candidates by a per-query
topic estimate and sets the number of documents and passages returned per
query by directly maximizing an expected-$F_2$ objective derived from the
reranker's own scores. A coverage-recovery step that resolves a subset of
the documents the primary crawl could not fetch directly, through their
archived snapshots, supplies additional documents to the same pipeline
without requiring a separate ranking mechanism.

On the public leaderboard of the task, the complete pipeline reaches a FINAL
score of 0.3060 (document-level $F_2$ 0.3381, passage-level $F_2$ 0.2739),
and each of the refinements described in Sect.~\ref{sec:fusion} --
topic-based re-ranking, per-query adaptive document counts, and per-query
passage windowing -- was associated with a further increase over the
configuration preceding it, for a combined gain of 0.0209 FINAL over a
fixed-size pseudo-relevance-feedback baseline. We also reported the design
alternatives that did not improve the score, including two -- candidate-pool
expansion by topic similarity and down-weighting the collection's dominant
language -- that a reader familiar only with the query language might expect
to help. The main limitation of this evaluation, that every comparison was
made directly on a rate-limited public leaderboard rather than on an
independent labelled set, is one we were unable to avoid under the task's
constraints, and we have tried to report it plainly rather than understate
it.

% ================= ACKNOWLEDGMENT & REFERENCES =================
\section*{Acknowledgment}
This research was supported by The VNUHCM-University of Information Technology's Scientific Research Support Fund.

\bibliographystyle{splncs04}
\bibliography{ref}

\end{document}
